\subsection{Sections and Element Properties}
\label{sec:sections}

Element connectivity defines interpolation and kinematics, but it does not by itself provide a complete structural model. Sections connect element regions to material and geometric properties. For solids the Section selects a material and optionally a material orientation. Shells additionally require thickness or a generalized shell stiffness. Trusses require cross-sectional area. Beams combine material data with cross-section constants and a local section direction.

A Section defined for an element set inside a Part belongs to that Part. If the Part is instantiated several times, every Instance carries the same Section assignment. Material, thickness, area, and profile data thereby remain properties of the physical component, while the Instance placement determines where that component is located.

\keywordnestedpage{PROFILE}

\cmd{PROFILE} defines the cross-section constants used by native FEMaster beam Sections. A Profile is referenced by name from \cmd{BEAM SECTION}. It stores area, bending second moments, torsional constant, optional product of inertia, shear-centre offsets, reference-line offsets, and optional effective transverse-shear areas.

The data order is
\[
A,\ I_y,\ I_z,\ J_t,\ I_{yz},\ e_y,\ e_z,\ \mathrm{ref}_y,\ \mathrm{ref}_z,\ A_{sy},\ A_{sz}.
\]
Only the first four values are required. Product of inertia and offsets default to zero. If \(A_{sy}\) or \(A_{sz}\) is omitted or non-positive, FEMaster uses \(5A/6\). The shear areas are used by shear-flexible beam formulations such as \val{B31}; \val{B33} does not use them.

FEMaster uses the convention
\[
I_{yz}=\int_A yz\,\mathrm dA
\]
without a leading minus sign. Some texts and tools define the product of inertia with the opposite sign, so imported Section properties can differ by this convention.

The offsets \(e_y,e_z\) describe the relative position used for the section/shear-centre definition, while \(\mathrm{ref}_y,\mathrm{ref}_z\) describe the location of the beam reference line relative to the section reference point. In unsymmetric or eccentric Sections these offsets determine the different points about which axial force, shear, bending, and torsion act.

Profile constants follow the unit system of the model. In a consistent N--mm model, \(A\), \(A_{sy}\), and \(A_{sz}\) are in mm\(^2\), \(I_y\), \(I_z\), \(J_t\), and \(I_{yz}\) are in mm\(^4\), and all offsets are in mm.

\begin{keywordparameters}
\key{PROFILE} / \key{NAME} & required & Unique Profile name referenced by \cmd{BEAM SECTION}. \\
\end{keywordparameters}

\begin{keyworddata}
\val{A} & Cross-sectional area. \\
\val{Iy, Iz} & Second moments of area about the local section axes. \\
\val{Jt} & Torsional constant. \\
\val{Iyz} & Product of inertia using \(I_{yz}=\int_A yz\,\mathrm dA\). \\
\val{ey, ez} & Optional section/shear-centre offsets; default zero. \\
\val{refy, refz} & Optional reference-line offsets; default zero. \\
\val{Asy, Asz} & Optional effective transverse-shear areas; default \(5A/6\). \\
\end{keyworddata}

The profile supplies section constants that a BEAM SECTION can reference by name. It does not by itself assign a section to any element.

\exampleheading

\begin{inputdeck}
*PROFILE, NAME=RECT_EQUIV
200., 6666.67, 1666.67, 4577.,
0., 0., 0., 0., 0., 166.667, 166.667.
\end{inputdeck}

\keywordnestedpage{SOLID SECTION}

\cmd{SOLID SECTION} assigns material properties to a solid element region. Native FEMaster requires \key{ELSET} and \key{MATERIAL}; \key{MAT} is accepted as an alias. \key{ORIENTATION} is optional and defines the local basis in which directional material properties are interpreted.

No thickness or area is required because a solid element represents its physical volume directly through the mesh geometry. The Section determines the constitutive law of the elements and the material coordinate system in which an anisotropic law is evaluated.

For isotropic elasticity, changing the material orientation has no effect on the constitutive response because the stiffness is direction independent. For orthotropic or fully anisotropic material data, the orientation is physically significant and must be chosen consistently with the material axes.

\cmd{SOLID SECTION} also accepts a positive area data line
when its ELSET contains only truss elements, including \val{T3}
and \val{T3D2}. This special case is equivalent to an explicit
\cmd{TRUSS SECTION}; for an ordinary solid region no area
data line is permitted.

A Section applies one material system and orientation to its target element set. Different material systems or orientations within one geometric component therefore require separate named element sets and Section assignments. A target set containing incompatible element families does not define a uniform Section interpretation.

\begin{keywordparameters}
\key{ELSET} & required & Solid element set receiving the Section. \\
\key{MATERIAL} / \key{MAT} & required & Material assigned to the region. \\
\key{ORIENTATION} & optional & Named material coordinate system. \\
\end{keywordparameters}

\paragraph*{Solid element section}

The first form assigns material and, where needed, material orientation to a continuum-element set. For a pure truss ELSET the same keyword takes an area data line instead.

\begin{inputdeck}
*SOLID SECTION, ELSET=SOLID,
 MATERIAL=STEEL,
 ORIENTATION=MAT_AXES
\end{inputdeck}

\paragraph*{Truss cross-sectional area}
\begin{inputdeck}
*SOLID SECTION, ELSET=BARS, MATERIAL=STEEL
100.
\end{inputdeck}

\paragraph*{Explicit truss section}
\begin{inputdeck}
*TRUSS SECTION, ELSET=BARS, MATERIAL=STEEL
100.
\end{inputdeck}

\keywordnestedpage{SHELL SECTION}

\cmd{SHELL SECTION} assigns thickness and constitutive properties to shell elements. Native FEMaster supports two forms:
\begin{itemize}
  \item \val{TYPE=INTEGRATED}, where a material law is integrated through the physical shell thickness;
  \item \val{TYPE=ABD}, where the generalized membrane--bending and transverse-shear stiffnesses are supplied directly.
\end{itemize}

For \val{TYPE=INTEGRATED}, \key{ELSET} identifies the shell region and \key{MATERIAL} identifies the constitutive law. The data line supplies the physical shell thickness \(t\). The shell uses this thickness to form membrane, bending, shear, mass, and stress-recovery quantities according to the shell formulation.

For a homogeneous linear elastic shell, the familiar generalized relation has the form
\[
\begin{bmatrix}
\boldsymbol N\\
\boldsymbol M
\end{bmatrix}
=
\begin{bmatrix}
A & B\\
B & D
\end{bmatrix}
\begin{bmatrix}
\boldsymbol\varepsilon_0\\
\boldsymbol\kappa
\end{bmatrix},
\]
with additional transverse-shear resultants governed by a \(2\times2\) shear stiffness. In an integrated Section these generalized stiffnesses arise from the material and thickness. In an ABD Section they are entered directly.

For \val{TYPE=ABD}, FEMaster reads 40 numerical values: all 36 entries of the full \(6\times6\) generalized stiffness matrix in row-major order, followed by the four entries of the \(2\times2\) transverse-shear stiffness matrix. The matrix is not assumed to be symmetric from input; enter all values explicitly in the required order.

\key{ORIENTATION} selects the local Section/material coordinate system. The ABD matrix is interpreted in this local shell basis. This is critical for laminates and anisotropic panels: an ABD matrix entered in one material direction but assigned with another orientation describes a different structure. \key{CSYSAXIS}=1,2,3 selects which axis of the named coordinate system establishes the in-plane reference direction; the default is axis 1.

The supported Abaqus \cmd{SHELL SECTION} form represents a homogeneous integrated Section. Its data line gives thickness and optionally the number of Simpson integration points. FEMaster currently supports five through-thickness points for this imported form; if the count is omitted, five are used. Layered composite layups and other Abaqus-specific shell-section options are not implied by this syntax.

The shell normal and in-plane Section direction together define the local shell basis in which \(N_{xx}\), \(N_{yy}\), \(M_{xx}\), \(M_{yy}\), top/bottom stresses, and an anisotropic ABD matrix are interpreted. On curved meshes both directions vary with the surface geometry. A 90-degree orientation difference can exchange strong and weak laminate directions while leaving the input syntactically valid.

\begin{keywordparameters}
\key{TYPE} & native optional & \val{INTEGRATED} (default) or \val{ABD}. \\
\key{ELSET} & required & Shell element set. \\
\key{MATERIAL} / \key{MAT} & integrated required & Material used by the integrated Section; optional for native ABD input. \\
\key{THICKNESS} & native optional & Nominal thickness parameter; default \val{1.0}. \\
\key{ORIENTATION} & optional & Named Section/material coordinate system. \\
\key{CSYSAXIS} & native optional & Coordinate-system axis 1, 2, or 3 used for the in-plane reference direction; default 1. \\
\end{keywordparameters}

\begin{keyworddata}
Integrated FEMaster & One positive shell thickness. \\
ABD FEMaster & 36 row-major generalized-stiffness entries followed by 4 row-major transverse-shear entries. \\
Abaqus homogeneous form & Thickness and optional integration-point count; the current supported count is 5. \\
\end{keyworddata}

\paragraph*{Integrated shell section}

An integrated shell section uses a material and thickness; a directly specified ABD section instead supplies membrane, bending, coupling, and transverse-shear stiffnesses.

\begin{inputdeck}
*SHELL SECTION, ELSET=PANEL,
 MATERIAL=ALUMINIUM
2.0, 5
\end{inputdeck}

\paragraph*{Direct ABD stiffness}
\begin{inputdeck}
*SHELL SECTION, ELSET=PANEL, TYPE=ABD
1000000., 0., 0., 0., 0., 0., 0., 1000000., 0., 0., 0., 0., 0., 0., 400000., 0., 0., 0., 0., 0., 0., 5000., 0., 0., 0., 0., 0., 0., 5000., 0., 0., 0., 0., 0., 0., 2000., 100000., 0., 0., 100000.
\end{inputdeck}

The data line contains 40 values: the complete $6\times6$
generalized membrane--bending stiffness matrix in row-major order,
followed by the complete $2\times2$ transverse-shear stiffness matrix.

\keywordnestedpage{BEAM SECTION}

\cmd{BEAM SECTION} and \cmd{BEAM GENERAL SECTION} both create the same generalized linear \val{BeamSection} in FEMaster. Unlike Abaqus, FEMaster does not yet integrate nonlinear constitutive behavior through a beam cross-section. Both input keywords therefore use the same parsing and section-property code path.

There are three mutually exclusive section descriptions:

\begin{itemize}
  \item \key{PROFILE=name} references a previously declared native FEMaster \cmd{PROFILE}; a named \key{MATERIAL} is required.
  \item \key{SECTION=GENERAL} (or an omitted \key{SECTION} without \key{PROFILE}) reads \((A,I_{11},I_{12},I_{22},J)\) directly on the first data line.
  \item \key{SECTION=RECT}, \val{CIRC}, \val{PIPE}, \val{BOX}, \val{I}, \val{CHANNEL}, or \val{L} reads a geometric data line and converts it to generalized section constants.
\end{itemize}

For geometric types, use Abaqus dimension order: \val{RECT}: \(a,b\); \val{CIRC}: radius; \val{PIPE}: outer radius, wall thickness; \val{BOX}: \(a,b,t\) for uniform walls or \(a,b,t_1,t_2,t_3,t_4\); \val{I}: \(l,h,b_1,b_2,t_1,t_2,t_3\), with a T-section represented by \(b_1=t_1=0\) or \(b_2=t_2=0\); \val{CHANNEL}: \(l,h,b_1,b_2,t_1,t_2,t_3,o\); \val{L}: \(a,b,t_1,t_2\). Shapes not listed here, including \val{ARBITRARY}, \val{HEX}, \val{HAT}, and \val{TRAPEZOID}, are not yet supported.

For \val{GENERAL} and geometric input, the first data line is followed by an optional \((n_{1x},n_{1y},n_{1z})\) direction (default \(0,0,-1\)). When \key{MATERIAL} is omitted, provide a following data line with \(E,G,\alpha\); \(\alpha\) is optional. The solver currently represents these direct constants using an equivalent isotropic material, so they must imply a valid Poisson ratio.

A \key{DENSITY} override may be specified without changing a shared named material. If \key{PROFILE} is used, the only data line is the optional section direction.

For closed circular sections, torsion is calculated from the polar second moment. For solid rectangles and thin-walled open/closed library shapes, torsion uses engineering approximations. The currently generic transverse shear areas remain \(5A/6\) and nontrivial shear-center corrections are not derived. Consequently, these inputs are linear generalized-section approximations, not a claim of full Abaqus library-section equivalence.

The nonzero orientation vector must not be parallel to the beam tangent. It defines the local cross-section basis in which section inertias and offsets are interpreted.

\begin{keywordparameters}
\key{ELSET} & required & Target beam element set. \\
\key{MATERIAL} / \key{MAT} & optional & Named material; required with \key{PROFILE}. \\
\key{PROFILE} & optional & Previously defined native Profile; mutually exclusive with \key{SECTION}. \\
\key{SECTION} & optional & \val{GENERAL} (default without PROFILE), or one of the supported library shapes. \\
\key{DENSITY} & optional & Section-local nonnegative density override. \\
\end{keywordparameters}

\paragraph*{Geometric beam section}

This beam section defines the geometry of a rectangular profile and its local reference direction. A separately defined PROFILE can be used when the section constants are already known.

\begin{inputdeck}
*BEAM SECTION, ELSET=BEAMS, MATERIAL=STEEL, SECTION=RECT
10., 20.
0., 1., 0.
\end{inputdeck}



The keyword \cmd{BEAM GENERAL SECTION} is also accepted with the
same \key{SECTION} geometry and data lines; only the keyword
spelling changes. No separate definition is required.

\paragraph*{Predefined profile}
\begin{inputdeck}
*PROFILE, NAME=BEAM_PROFILE
200., 6666.7, 1666.7, 4577.
*BEAM SECTION, ELSET=BEAMS, MATERIAL=STEEL, PROFILE=BEAM_PROFILE
0., 1., 0.
\end{inputdeck}

\keywordnestedpage{TRUSS SECTION}

Native \cmd{TRUSS SECTION} assigns a material and a positive cross-sectional area to a truss element set. A truss carries only axial force, so cross-sectional area is the only geometric Section property required by the current formulation.

For an element of length \(L\), area \(A\), and Young's modulus \(E\), the small-strain axial stiffness scales as
\[
k=\frac{EA}{L}.
\]
The same area also determines material mass per unit length through \(\rho A\) and converts axial force to average axial stress through \(\sigma=N/A\).

The area must therefore use the same unit system as the geometry and material constants and must be strictly positive.

\cmd{TRUSS SECTION} is the explicit FEMaster form for assigning
a truss material and area. Alternatively, \cmd{SOLID SECTION}
accepts an area data line when the ELSET consists exclusively
of truss elements (\val{T3} or \val{T3D2}). The resulting
truss-section properties are the same. Ordinary solid elements,
in contrast, use \cmd{SOLID SECTION} without an area line.

\begin{keywordparameters}
\key{MATERIAL} / \key{MAT} & required & Truss material. \\
\key{ELSET} & required & Truss element set. \\
\end{keywordparameters}

\begin{keyworddata}
\val{area} & Strictly positive cross-sectional area. \\
\end{keyworddata}

A truss carries axial force only. The area given on the data line is the area used for axial stiffness and mass; SOLID SECTION can assign the same area to an exclusively truss-based set.

\exampleheading

\begin{inputdeck}
*TRUSS SECTION, ELSET=BARS,
 MATERIAL=STEEL
100.0
\end{inputdeck}

\paragraph*{Solid section for truss elements}
\begin{inputdeck}
*SOLID SECTION, ELSET=BARS,
 MATERIAL=STEEL
100.0
\end{inputdeck}
